Artificial Intelligence (AI) techniques in games have been developed for decades, progressing the field of computer-aided decision-making with them. Game tree search is a particular branch of game AI. This discipline represents the game as a rooted tree, for which the best node at its root is the next move to play. B* is a best-first game-tree search technique developed in the late 1970s, which was later used for competitive endgame play of Scrabble in the 2000s. B* is unique in its construction and capabilities, but has not been extensively researched in the last two decades. A recent report has shown that B* can solve large problems more efficiently than was previously known. However, B* is limited in the size of problems it can solve, due to its memory requirements. In this thesis, we address this problem by introducing two B* variants: the memory-efficient multi-level B*² and the memory-independent depth-first df-B*. 
Experimental evaluation on over 37,000 artificial game trees showed B*² requires just 7.6 times the square root of the memory usage of B*. Whereas the experiments showed that df-B* only requires a near-constant amount of memory, effectively never running out of memory on any problem. However, we observed that df-B* has significantly reduced computational efficiency by approximately a factor of 22 to 23 compared to B*. By contrast, we found B*² was able to solve some problems more efficiently than B*, and was, on average, only a factor of up to 3 times less computationally efficient. In conclusion, B*² was found to be effective for search on large problems, whereas df-B* was found to be effective for search on very large problems, being effectively unrestrained by memory limitations. These results were found to align with existing research on memory-efficient variants of the best-first game-tree solver Proof-Number Search.
\\
\phantom{}
\vspace{2ex}\\
\textbf{Keywords:} B* Search $\cdot$ Artificial Game Trees $\cdot$  Depth-First Search $\cdot$ Adversarial Search $\cdot$ Artificial Intelligence